\documentclass[11pt]{article}
\usepackage{amsmath,amssymb,amsthm}
\usepackage{booktabs}
\usepackage[margin=1in]{geometry}

\newtheorem{theorem}{Theorem}
\newtheorem{lemma}[theorem]{Lemma}
\newtheorem{corollary}[theorem]{Corollary}
\newtheorem{proposition}[theorem]{Proposition}

\title{Problem 625: GCDSUM}
\author{Project Euler Solutions}
\date{}

\begin{document}
\maketitle

\section{Problem Statement}

Define $G(N) = \sum_{j=1}^{N}\sum_{i=1}^{j}\gcd(i,j)$. Compute $G(10^{11}) \bmod 998244353$.

\section{Totient Identity}

\begin{theorem}[Gauss]
For every positive integer~$n$:
\[
n = \sum_{d \mid n} \varphi(d).
\]
\end{theorem}

\begin{proof}
Partition $\{1, \ldots, n\}$ by $\gcd(k, n)$: for each $d \mid n$, exactly $\varphi(n/d)$ elements satisfy $\gcd(k,n) = d$. Equivalently, $\sum_{d \mid n}\varphi(d) = n$.
\end{proof}

\begin{corollary}
$\gcd(i,j) = \sum_{d \mid \gcd(i,j)} \varphi(d)$.
\end{corollary}

\section{Sum Transformation}

\begin{theorem}
\[
G(N) = \sum_{d=1}^{N} \varphi(d) \cdot T\!\left(\left\lfloor\frac{N}{d}\right\rfloor\right)
\]
where $T(m) = m(m+1)/2$.
\end{theorem}

\begin{proof}
Replace $\gcd$ by totient and swap summation:
\[
G(N) = \sum_{d=1}^N \varphi(d) \!\!\sum_{\substack{1 \le i \le j \le N \\ d \mid i,\; d \mid j}} \!\!1 = \sum_{d=1}^N \varphi(d) \binom{\lfloor N/d\rfloor + 1}{2}. \qedhere
\]
\end{proof}

\section{Block Decomposition}

\begin{proposition}
$\lfloor N/d \rfloor$ takes at most $2\sqrt{N}$ distinct values as $d$ ranges from~$1$ to~$N$.
\end{proposition}

For each block $[l, r]$ where $q = \lfloor N/d \rfloor$ is constant:
\[
\sum_{d=l}^{r} \varphi(d)\,T(q) = T(q)\bigl(\Phi(r) - \Phi(l-1)\bigr)
\]
where $\Phi(n) = \sum_{k=1}^n \varphi(k)$ is the \textbf{totient summatory function}.

\section{Sub-Linear Totient Summation}

\begin{lemma}
\[
\Phi(n) = \frac{n(n+1)}{2} - \sum_{d=2}^{n} \Phi\!\left(\left\lfloor\frac{n}{d}\right\rfloor\right).
\]
\end{lemma}

\begin{proof}
From $\sum_{n=1}^N \sum_{d \mid n}\varphi(d) = \sum_{n=1}^N n = N(N+1)/2$. The left side equals $\sum_{d=1}^N \varphi(d)\lfloor N/d\rfloor = \Phi(N) + \sum_{d=2}^N \Phi(\lfloor N/d\rfloor)$ after rearranging.
\end{proof}

This identity, applied recursively with memoization and block decomposition internally, computes $\Phi(n)$ for all needed arguments in $O(N^{2/3})$ total time.

\section{Concrete Values}

\begin{center}
\begin{tabular}{ccc}
\toprule
$N$ & $G(N)$ & $G(N) / N^2$ \\
\midrule
1   & 1      & 1.000 \\
2   & 4      & 1.000 \\
5   & 29     & 1.160 \\
10  & 138    & 1.380 \\
100 & 12614  & 1.261 \\
\bottomrule
\end{tabular}
\end{center}

\section{Complexity Analysis}

\begin{itemize}
\item \textbf{Sieve:} $O(N^{2/3})$ for Euler's totient up to $N^{2/3}$.
\item \textbf{$\Phi$ computation:} $O(N^{2/3})$ total over all recursive calls.
\item \textbf{Block sum:} $O(\sqrt{N})$ blocks for the outer sum.
\item \textbf{Total:} $O(N^{2/3})$ time, $O(N^{2/3})$ space.
\end{itemize}

For $N = 10^{11}$, $N^{2/3} \approx 2.15 \times 10^7$.

\section{Answer}

\[
\boxed{551614306}
\]

\end{document}
